\chapter{Metodología y planificación}
\label{chap:metodologia}

\section{Enfoque de trabajo}
\label{sec:enfoque-trabajo}

El proyecto se desarrolla con un enfoque iterativo e incremental. Cada iteración busca producir una versión funcional o verificable del sistema, evitando introducir complejidad que no aporte valor directo a los objetivos del trabajo.

\section{Fases del proyecto}
\label{sec:fases-proyecto}

\begin{enumerate}
  \item Análisis y diseño del alcance, requisitos, modelo de datos y arquitectura.
  \item Desarrollo del backend y de los casos de uso principales.
  \item Contenerización y preparación del entorno local reproducible.
  \item Definición de la infraestructura cloud mediante Terraform.
  \item Automatización de integración y despliegue continuo.
  \item Incorporación de observabilidad y medidas básicas de seguridad.
  \item Validación funcional y análisis del comportamiento en ejecución.
  \item Redacción de memoria, conclusiones y líneas futuras.
\end{enumerate}

\section{Gestión del alcance}
\label{sec:gestion-alcance}

La gestión del alcance se ha basado en una decisión explícita: construir un producto funcionalmente pequeño, pero técnicamente completo. En vez de añadir módulos de negocio avanzados, como optimización de rutas, facturación, geolocalización en tiempo real o panel web, se priorizó que el backend mínimo pudiera desplegarse y operarse en AWS de forma creíble.

Esta decisión reduce el riesgo de terminar con una aplicación amplia pero superficial. El valor del trabajo se concentra en demostrar el ciclo completo de ingeniería: modelado de dominio, API, persistencia, pruebas, contenedores, infraestructura cloud, seguridad básica, CI/CD, migraciones y validación operativa.

El alcance funcional se fijó alrededor de transportes, vehículos, conductores, usuarios, asignaciones, estados y eventos logísticos. Cualquier funcionalidad fuera de ese flujo se trató como línea futura salvo que fuese necesaria para demostrar el despliegue o la operación.

\section{Herramientas utilizadas}
\label{sec:herramientas}

Las herramientas principales utilizadas han sido:

\begin{itemize}
  \item ASP.NET Core y .NET 10 para la implementación de la API REST~\cite{MicrosoftAspNetCoreDocs}.
  \item Entity Framework Core como ORM y herramienta de migraciones.
  \item PostgreSQL como base de datos relacional.
  \item xUnit y \texttt{WebApplicationFactory} para pruebas automatizadas.
  \item Docker y Docker Compose para empaquetado y ejecución local reproducible.
  \item Terraform para definir infraestructura AWS como código.
  \item AWS ECS Fargate, ECR, RDS, ALB, Route 53, ACM, CloudWatch, Secrets Manager y SSM Parameter Store como servicios cloud principales~\cite{AwsEcsDocs,AwsRdsDocs,AwsRoute53Docs,AwsAcmDocs}.
  \item GitHub Actions para integración continua y despliegue manual.
  \item Postman/Newman para pruebas de humo manuales y automatizables.
  \item LaTeX para la elaboración de la memoria.
\end{itemize}
